\section{Conditional Trajectory Abstraction}
\label{sec:method}

\method is a joint-embedding predictive world model.
A target encoder turns the actual future into a code (\cref{sec:source}), two heads read the code (\cref{sec:reader}), and a predictor forecasts the code from the action (\cref{sec:wm}).
\Cref{sec:codesign} describes training, \cref{sec:planning} planning, and \cref{sec:impl} implementation details.

\subsection{Inputs}
\label{sec:context}
Images are embedded by a frozen DINOv2 ViT-S/14~\cite{oquab2024dinov2}, following DINO-WM~\cite{zhou2025dinowm}.
We keep the full $16{\times}16$ patch grid and project the patch features to 128 dimensions with a PCA fitted on training frames.
The context $C$ contains the patch tokens of the current image, the tokens of the previous image pooled to an $8{\times}8$ grid, and the proprioceptive states at both times.
The future segment $\tau$ contains the tokens of the segment's final image on the full grid, the tokens of intermediate images pooled to $8{\times}8$ with temporal embeddings, and the final proprioceptive state.
Features are computed once per image and shared by all candidates and goals.
Each module embeds $C$ and $\tau$ with its own linear projections and learned position and type embeddings.

\subsection{Target encoder}
\label{sec:source}
The target encoder $q_\phi(S\mid C,\tau)$ exists only during training.
$M$ learned segment queries cross-attend jointly to the future tokens and the context tokens in a Transformer decoder.
Each output is projected to three dimensions and quantized by finite scalar quantization (FSQ)~\cite{mentzer2024fsq} with levels $(8,8,4)$, an implicit vocabulary of $V{=}256$ per token.
Because the queries attend to the context, the encoder can learn which parts of the future the history already implies, instead of relying on a hand-designed residual such as the difference between future and current features.
The code describes the whole segment; it is not a set of $M$ tokens per time step.

\subsection{Reader and feature decoder}
\label{sec:reader}
\paragraph{Reader.}
The reader $D_\psi(C,S,g)$ is a Transformer over a summary token, the context tokens, the embedded code tokens, and the patch tokens of a goal image $g$.
It outputs a scalar score from the summary token.
It never receives the action, so the code is the only channel through which the action's consequences reach it.
The reader is trained on the $K$ candidates branched from one state, whose futures and labels are known during training.
It uses a pairwise logistic ranking loss over the pairs whose labels differ by more than a margin $\epsilon$:
\begin{equation}
\begin{aligned}
  \mathcal{L}_{\mathrm{rank}}&=\E_{(i,j):\,\ell_i>\ell_j+\epsilon}\big[\log\big(1+e^{-(s_i-s_j)}\big)\big],\\
  s_i&=D_\psi(C,S_i,g),\quad \ell_i=\ell_g(\tau_i).
\end{aligned}
\label{eq:rank}
\end{equation}
Ranking within a bank matches how the score is used: it compares candidates from one state and need not be calibrated across states.
These labels are the only task supervision in \method, and they are used in training only.

\paragraph{Feature decoder.}
A ranking loss alone lets the code shrink to whatever orders the training goals, which would not transfer to new goals.
The feature decoder $G_\psi(C,S)$ therefore reconstructs the future tokens from the context and the code.
One learned query per future token position cross-attends to the context and code tokens, and the decoder predicts how each token changes from the corresponding token of the current image.
Its loss $\mathcal{L}_{\mathrm{feat}}$ is the squared error divided by that of a copy baseline that predicts no change, computed separately for the final and the intermediate images, so a code that carries nothing about the change scores one.
In JEPA terms, the feature decoder is the regularizer that keeps the target representation from collapsing~\cite{lecun2022path}.
The world model's prediction target remains the code, and no image is ever generated.

\subsection{Action-conditioned world model}
\label{sec:wm}
The world model $p_\theta(S\mid C,\bA)$ is the predictor of the joint-embedding architecture.
A Transformer encoder attends over the context tokens and one token per executed action.
A Transformer decoder then generates the $M$ code tokens autoregressively as categorical distributions over the vocabulary, which preserves the dependencies between code positions.
The world model never receives future images, future proprioception, the goal, or any label.
It is trained by teacher forcing against stop-gradient codes of the target encoder,
\begin{equation}
  \mathcal{L}_{\mathrm{WM}}=\E\big[-\log p_\theta(\sg(S)\mid C,\bA)\big].
  \label{eq:wm}
\end{equation}
A context-only model $r_\omega(S\mid C)$ with the same architecture but no action tokens is trained alongside.
It estimates the rate of the code.
Its excess cross-entropy over the world model estimates $I(S;\bA\mid C)$, the part of the code that the action explains; a code that ignores the action would show none.
At test time we decode the most likely code greedily.
We also evaluate two variants: the expected FSQ codeword under the world model at each position, given the greedy prefix, and the reader's scores averaged over four sampled codes.

\subsection{Training}
\label{sec:codesign}
\paragraph{Data.}
Training tuples come from rollouts of the frozen proposal policy.
At every decision of a rollout, all $K$ proposals are executed for $\He$ steps from clones of the simulator state.
This yields $K$ futures and their labels from an identical history, and the rollout then continues with the default proposal.
All branches of a root episode belong to the same data split.
Branching is what lets both the ranking loss and the world model observe different actions from the same history.

\paragraph{Stage 1: code.}
The target encoder, reader, and feature decoder are trained jointly on actual futures with
\begin{equation}
  \mathcal{L}_{\mathrm{code}}=\mathcal{L}_{\mathrm{rank}}+\alpha\,\mathcal{L}_{\mathrm{feat}}.
  \label{eq:code}
\end{equation}
One goal image per decision is sampled from the training goals, and gradients pass through the quantizer with the straight-through estimator.

\paragraph{Stage 2: world model.}
With the target encoder frozen, the world model and the context-only model are trained on its codes by~\eqref{eq:wm}.

\paragraph{Stage 3: co-design (optional).}
A code that retains the future well may still be hard to forecast.
In this stage, fixed blocks of code updates alternate with fixed blocks of world-model updates.
Code updates add a predictability term under a frozen copy $\bar p_\theta$ of the world model:
\begin{equation}
\begin{aligned}
  \mathcal{L}^{\lambda}_{\mathrm{code}}={}&\mathcal{L}_{\mathrm{rank}}+\alpha\,\mathcal{L}_{\mathrm{feat}}\\
  &+\lambda\,\big\lVert S-\E_{\bar p_\theta(\cdot\mid C,\bA)}[S]\big\rVert^2.
\end{aligned}
  \label{eq:source}
\end{equation}
The last term is a differentiable surrogate of $-\log\bar p_\theta(S\mid C,\bA)$.
It pulls each quantized token, viewed as its FSQ codeword, toward the codeword the frozen world model expects given the preceding tokens.
World-model updates use targets from an exponential moving average of the target encoder, which replaces the encoder at the end.
Setting $\lambda{=}0$ reduces training to stages 1 and 2, and this ablation tests whether predictability-aware co-design adds anything.
\TODO{Optional reader adaptation (fine-tuning $D_\psi$ on a mix of target and predicted codes with everything else frozen) is a debugging option; describe it here only if the locked configuration uses it.}

\paragraph{Collapse diagnostics.}
A code can satisfy the losses while ignoring the action or the future.
We monitor code perplexity, the fraction of distinct codes within a bank, the feature decoder's $R^2$ relative to the copy baseline, and the action information $I(S;\bA\mid C)$.

\subsection{Planning with predicted codes}
\label{sec:planning}
\Cref{alg:cta} summarizes deployment.
The planning score of a candidate is the reader's score averaged over the goal images, which are final frames of successful episodes of the task.
All candidates are scored from the context and their actions, before any of them is executed.
Per decision, \method computes the image features once, runs the world model on the $K$ candidates in one batch, and runs the reader on $K|\cG|$ candidate--goal pairs.
Because codes do not depend on the goal, a new goal on the same candidates costs only additional reader evaluations.

\begin{algorithm}[t]
\caption{Policy-guided planning with \method.}
\label{alg:cta}
\small
\begin{algorithmic}[1]
\Require frozen policy $\pi$; frozen encoder $\varphi$; world model $p_\theta$; reader $D_\psi$; goal images $\cG$; $K$; $\He$
\While{episode not finished}
  \State $C\gets$ features of $h_t$ under $\varphi$ \Comment{once per decision}
  \State $\bA^{(k)}\gets\pi(h_t;\xi_k)$ for $k=1,\dots,K$ \Comment{shared bank}
  \State $\hat S^{(k)}\gets$ greedy decode of $p_\theta(\cdot\mid C,\bA^{(k)})$ \Comment{batched}
  \State $s_k\gets\frac{1}{|\cG|}\sum_{g\in\cG}D_\psi\big(C,\hat S^{(k)},g\big)$
  \State $k^\star\gets$ smallest index attaining $\max_k s_k$ \Comment{ties keep the default}
  \State execute $\bA^{(k^\star)}_{1:\He}$; $t\gets t+\He$
\EndWhile
\end{algorithmic}
\end{algorithm}

\subsection{Implementation details}
\label{sec:impl}
All modules have width 256 and 8 attention heads.
The target encoder has three decoder layers, the reader four encoder layers, and the feature decoder two decoder layers.
The world model and the context-only model each have four encoder and four decoder layers.
Unless stated otherwise, $M{=}16$ (128 bits per candidate), $\alpha{=}1$, and $\epsilon{=}10^{-3}$.
We train with AdamW (learning rate $3{\times}10^{-4}$, weight decay $0.05$, 500 warm-up steps, cosine decay) in bfloat16.
Stage 1 runs 10k steps with 16 decisions (128 candidates) per batch, and stage 2 runs 20k steps with 32 decisions per batch.
Stage 3 alternates 200 code steps with 400 world-model steps at learning rate $10^{-4}$, with EMA rate $0.99$.
On PushT, images are $96{\times}96$ and resized to $224{\times}224$ for DINOv2, $\He{=}8$, the intermediate images are those after steps 2, 4, and 6, and $K{=}8$.
Each action is a target position of the agent, embedded both in absolute coordinates and relative to the current position.
\TODO{LIBERO-Goal: action embedding, $\He$, intermediate frames, camera views.}
